% Ported to the default house style by the handout-formatting skill, from
% inputs/math/Worksheets/03-First-Order-Linear-Differential-Equations.pdf.
% One source, three audiences (change only this option line):
%   \usepackage[]{handout}        student handout (blanks, no solutions)
%   \usepackage[key]{handout}     answer key (solutions shown)
%   \usepackage[teacher]{handout} teacher copy (solutions + teacher notes)
\documentclass[10pt]{article}
\usepackage{housestyle}
\usepackage[teacher]{handout}

\hypersetup{pdftitle={First Order Linear Differential Equations -- Worksheet 3}}

\begin{document}

\handouttitle{First Order Linear Differential Equations}{Worksheet 3}{}

\calloutbox{Today's Goals}{%
    \begin{itemize}[leftmargin=*]
        \item Recognize a first order \textbf{linear} equation \(\frac{dy}{dx} + P(x)y = Q(x)\).
        \item Compute the integrating factor \(\rho(x) = e^{\int P(x)\,dx}\).
        \item Use \(\frac{d}{dx}[\rho(x)y] = \rho(x)Q(x)\) to integrate both sides.
        \item Solve general and initial-value problems, including ones requiring substitution to integrate.
    \end{itemize}%
}

\bigskip

We solve linear first order equations of the form
\[
    \frac{dy}{dx} + P(x)y = Q(x),
\]
where \(P(x)\) and \(Q(x)\) are continuous. We cannot separate variables, but multiplying by the
\textbf{integrating factor} \(\rho(x) = e^{\int P(x)\,dx}\) makes the left side an exact derivative.

\begin{note}
    Since \(\dfrac{d}{dx}\!\left(\int P(x)\,dx\right) = P(x)\), the product rule gives
    \[
        \frac{d}{dx}\!\left[ y\, e^{\int P(x)\,dx}\right]
        = e^{\int P(x)\,dx}\frac{dy}{dx} + P(x)e^{\int P(x)\,dx}\,y
        = Q(x)\, e^{\int P(x)\,dx}.
    \]
    Integrating, \(\ y\,e^{\int P(x)\,dx} = \int Q(x)\, e^{\int P(x)\,dx}\,dx + C\).
\end{note}

\textbf{Steps to solve} \(\dfrac{dy}{dx} + P(x)y = Q(x)\):
\begin{enumerate}[leftmargin=*, label=\arabic*.]
    \item Calculate \(\rho(x) = e^{\int P(x)\,dx}\).
    \item Multiply the equation by \(\rho(x)\).
    \item Notice that \(\dfrac{d}{dx}\!\left[\rho(x)y\right] = \rho(x)Q(x)\).
    \item Integrate both sides: \(\rho(x)y = \int \rho(x)Q(x)\,dx + C\).
    \item Solve: \(y = \dfrac{1}{\rho(x)}\!\left[\int \rho(x)Q(x)\,dx + C\right]\).
\end{enumerate}

\section{Using an Integrating Factor}

\begin{enumerate}[leftmargin=*, label=\numbox{\arabic*}]

\item
    \begin{problem}[6cm]
        Solve \(\ 2xy' + y = 10\sqrt{x}\ \) for \(x > 0\).
    \end{problem}
    \begin{solution}
        Put the equation in standard form by dividing by \(2x\):
        \[
            y' + \frac{1}{2x}\,y = \frac{5}{\sqrt{x}},
            \qquad P(x) = \frac{1}{2x}, \quad Q(x) = \frac{5}{\sqrt{x}}.
        \]
        The integrating factor is
        \[
            \rho(x) = e^{\int \frac{1}{2x}\,dx} = e^{\frac{1}{2}\ln x} = (e^{\ln x})^{1/2} = x^{1/2}.
        \]
        Multiplying by \(\sqrt{x}\):
        \[
            \sqrt{x}\,y' + \frac{1}{2\sqrt{x}}\,y = 5,
            \qquad\text{i.e.}\qquad \frac{d}{dx}\!\left(\sqrt{x}\,y\right) = 5.
        \]
        Integrating, \(\ \sqrt{x}\,y = 5x + C\), so
        \[
            y = \frac{1}{\sqrt{x}}(5x + C) = 5\sqrt{x} + \frac{C}{\sqrt{x}} \qquad\text{(general solution).}
        \]
    \end{solution}

\item
    \begin{problem}[6cm]
        Solve the initial value problem \(\ y' = 2xy + 3x^2 e^{x^2}\), \(\ y(0) = 5\).
    \end{problem}
    \begin{solution}
        In standard form, \(\ y' - 2xy = 3x^2 e^{x^2}\), so \(P(x) = -2x\) and \(Q(x) = 3x^2 e^{x^2}\).
        \[
            \rho(x) = e^{\int -2x\,dx} = e^{-x^2}.
        \]
        Multiplying through:
        \[
            e^{-x^2}y' - 2x e^{-x^2}y = 3x^2 e^{x^2}\cdot e^{-x^2} = 3x^2,
            \qquad \frac{d}{dx}\!\left(e^{-x^2}y\right) = 3x^2.
        \]
        Integrating, \(\ e^{-x^2}y = x^3 + C\), so \(y = e^{x^2}(x^3 + C)\) (general solution).
        Applying \(y(0) = 5\):
        \[
            5 = e^{0}(0 + C) = C,
        \]
        so \(\ y = e^{x^2}(x^3 + 5)\) (particular solution).
    \end{solution}

\item
    \begin{problem}[7cm]
        Solve \(\ y' = -y\tan x + (\cos^2 x)\sin x\), \(\ y(0) = 3\), for \(-\tfrac{\pi}{2} < x < \tfrac{\pi}{2}\).
    \end{problem}
    \begin{solution}
        In standard form, \(\ y' + (\tan x)y = (\cos^2 x)\sin x\), so \(P(x) = \tan x\) and
        \(Q(x) = (\cos^2 x)\sin x\). With \(u = \cos x\), \(du = -\sin x\,dx\):
        \[
            \rho(x) = e^{\int \tan x\,dx} = e^{\int \frac{\sin x}{\cos x}\,dx}
            = e^{-\int \frac{du}{u}} = e^{-\ln u} = \frac{1}{\cos x} = \sec x,
        \]
        valid since \(\cos x > 0\) on \(-\tfrac{\pi}{2} < x < \tfrac{\pi}{2}\). Multiplying through:
        \[
            (\sec x)y' + (\sec x)(\tan x)y = (\sec x)(\cos^2 x)\sin x = (\cos x)\sin x,
        \]
        i.e. \(\dfrac{d}{dx}\!\left[(\sec x)y\right] = (\cos x)\sin x\). Integrating with \(u = \cos x\):
        \[
            (\sec x)y = \int (\cos x)\sin x\,dx = -\int u\,du = -\frac{u^2}{2} + C = -\frac{\cos^2 x}{2} + C.
        \]
        Thus \(\ y = -\dfrac{\cos^3 x}{2} + C\cos x\) (general solution). Applying \(y(0) = 3\):
        \[
            3 = -\frac{\cos^2 0}{2} + C\cos 0 = -\frac{1}{2} + C \implies C = \frac{7}{2},
        \]
        so \(\ y = -\dfrac{\cos^3 x}{2} + \dfrac{7}{2}\cos x\) (particular solution).
    \end{solution}

\end{enumerate}

\end{document}
